\begin{register}{H}{LEDC\_HSCH\textcolor{red}{\it{n}}\_CONF0\_REG (\textcolor{red}{\it{n}}: 0-7)}{0x00+0x14*\textcolor{red}{\it{n}}}\label{regdesc:LEDCHSCHnCONF0REG}
\regfield{(reserved)}{28}{4}{{0x00000000}}%
\regfield{LEDC\_IDLE\_LV\_HSCH\textcolor{red}{\it{n}}}{1}{3}{0}%
\regfield{LEDC\_SIG\_OUT\_EN\_HSCH\textcolor{red}{\it{n}}}{1}{2}{0}%
\regfield{LEDC\_TIMER\_SEL\_HSCH\textcolor{red}{\it{n}}}{2}{0}{{0}}%
\reglabel{Reset}\regnewline%
\begin{regdesc}\begin{reglist}
\item [LEDC\_IDLE\_LV\_HSCH\textcolor{red}{\it{n}}]This bit is used to control the output value when high-speed channel \textcolor{red}{\it{n}} is inactive. (R/W)
\item [LEDC\_SIG\_OUT\_EN\_HSCH\textcolor{red}{\it{n}}]This is the output enable control bit for high-speed channel \textcolor{red}{\it{n}}. (R/W)
\item [LEDC\_TIMER\_SEL\_HSCH\textcolor{red}{\it{n}}]There are four high-speed timers. These two bits are used to select one of them for high-speed channel \textcolor{red}{\it{n}}:  (R/W) \\0: select hstimer0; \\1: select hstimer1;  \\2: select hstimer2; \\3: select hstimer3.
\end{reglist}\end{regdesc}
\end{register}


\begin{register}{H}{LEDC\_HSCH\textcolor{red}{\it{n}}\_HPOINT\_REG (\textcolor{red}{\it{n}}: 0-7)}{0x04+0x14*\textcolor{red}{\it{n}}}\label{regdesc:LEDCHSCHnHPOINTREG}
\regfield{(reserved)}{12}{20}{{0x0000}}%
\regfield{LEDC\_HPOINT\_HSCH\textcolor{red}{\it{n}}}{20}{0}{{0x000000}}%
\reglabel{Reset}\regnewline%
\begin{regdesc}\begin{reglist}
\item [LEDC\_HPOINT\_HSCH\textcolor{red}{\it{n}}]The output value changes to high when htimerx(x=[0,3]), selected by high-speed channel \textcolor{red}{\it{n}}, has reached LEDC\_HPOINT\_HSCH\textcolor{red}{\it{n}}[19:0]. (R/W)
\end{reglist}\end{regdesc}
\end{register}


\begin{register}{H}{LEDC\_HSCH\textcolor{red}{\it{n}}\_DUTY\_REG (\textcolor{red}{\it{n}}: 0-7)}{0x08+0x14*\textcolor{red}{\it{n}}}\label{regdesc:LEDCHSCHnDUTYREG}
\regfield{(reserved)}{7}{25}{{0x00}}%
\regfield{LEDC\_DUTY\_HSCH\textcolor{red}{\it{n}}}{25}{0}{{0x0000000}}%
\reglabel{Reset}\regnewline%
\begin{regdesc}\begin{reglist}
\item [LEDC\_DUTY\_HSCH\textcolor{red}{\it{n}}]The register is used to control output duty. When hstimerx(x=[0,3]), selected by high-speed channel \textcolor{red}{\it{n}}, has reached LEDC\_LPOINT\_HSCH\textcolor{red}{\it{n}}, the output signal changes to low. (R/W) \\ LEDC\_LPOINT\_HSCH\textcolor{red}{\it{n}}=LEDC\_LPOINT\_HSCH\textcolor{red}{\it{n}}[19:0]+LEDC\_DUTY\_HSCH\textcolor{red}{\it{n}}[24:4] (1) \\LEDC\_LPOINT\_HSCH\textcolor{red}{\it{n}}=LEDC\_LPOINT\_HSCH\textcolor{red}{\it{n}}[19:0]+LEDC\_DUTY\_HSCH\textcolor{red}{\it{n}}[24:4] +1) (2) \\  See the \hyperref[ledcfd:dithering]{Functional Description} for more information on when (1) or (2) is chosen.
\end{reglist}\end{regdesc}
\end{register}


\begin{register}{H}{LEDC\_HSCH\textcolor{red}{\it{n}}\_CONF1\_REG (\textcolor{red}{\it{n}}: 0-7)}{0x0C+0x14*\textcolor{red}{\it{n}}}\label{regdesc:LEDCHSCHnCONF1REG}
\regfield{LEDC\_DUTY\_START\_HSCH\textcolor{red}{\it{n}}}{1}{31}{0}%
\regfield{LEDC\_DUTY\_INC\_HSCH\textcolor{red}{\it{n}}}{1}{30}{1}%
\regfield{LEDC\_DUTY\_NUM\_HSCH\textcolor{red}{\it{n}}}{10}{20}{{0x000}}%
\regfield{LEDC\_DUTY\_CYCLE\_HSCH\textcolor{red}{\it{n}}}{10}{10}{{0x000}}%
\regfield{LEDC\_DUTY\_SCALE\_HSCH\textcolor{red}{\it{n}}}{10}{0}{{0x000}}%
\reglabel{Reset}\regnewline%
\begin{regdesc}\begin{reglist}
\item [LEDC\_DUTY\_START\_HSCH\textcolor{red}{\it{n}}]When LEDC\_DUTY\_NUM\_HSCH\textcolor{red}{\it{n}}, LEDC\_DUTY\_CYCLE\_HSCH\textcolor{red}{\it{n}} and LEDC\_DUTY\_SCALE\_HSCH\textcolor{red}{\it{n}} has been configured, these register will not take effect until LEDC\_DUTY\_START\_HSCH\textcolor{red}{\it{n}} is set. This bit is automatically cleared by hardware. (R/W)
\item [LEDC\_DUTY\_INC\_HSCH\textcolor{red}{\it{n}}]This register is used to increase or decrease the duty of output signal for high-speed channel \textcolor{red}{\it{n}}. (R/W)
\item [LEDC\_DUTY\_NUM\_HSCH\textcolor{red}{\it{n}}]This register is used to control the number of times the duty cycle is increased or decreased for high-speed channel \textcolor{red}{\it{n}}. (R/W)
\item [LEDC\_DUTY\_CYCLE\_HSCH\textcolor{red}{\it{n}}]This register is used to increase or decrease the duty cycle every time LEDC\_DUTY\_CYCLE\_HSCH\textcolor{red}{\it{n}} cycles for high-speed channel \textcolor{red}{\it{n}}. (R/W)
\item [LEDC\_DUTY\_SCALE\_HSCH\textcolor{red}{\it{n}}]This register is used to increase or decrease the step scale for high-speed channel \textcolor{red}{\it{n}}. (R/W)
\end{reglist}\end{regdesc}
\end{register}


\begin{register}{H}{LEDC\_HSCH\textcolor{red}{\it{n}}\_DUTY\_R\_REG (\textcolor{red}{\it{n}}: 0-7)}{0x10+0x14*\textcolor{red}{\it{n}}}\label{regdesc:LEDCHSCHnDUTYRREG}
\regfield{(reserved)}{7}{25}{{0x00}}%
\regfield{LEDC\_DUTY\_HSCH\textcolor{red}{\it{n}}\_R}{25}{0}{{0x0000000}}%
\reglabel{Reset}\regnewline%
\begin{regdesc}\begin{reglist}
\item [LEDC\_DUTY\_HSCH\textcolor{red}{\it{n}}\_R]This register represents the current duty cycle of the output signal for high-speed channel \textcolor{red}{\it{n}}. (RO)
\end{reglist}\end{regdesc}
\end{register}


\begin{register}{H}{LEDC\_LSCH\textcolor{red}{\it{n}}\_CONF0\_REG (\textcolor{red}{\it{n}}: 0-7)}{0xA0+0x14*\textcolor{red}{\it{n}}}\label{regdesc:LEDCLSCHnCONF0REG}
\regfield{(reserved)}{27}{5}{{0x0000000}}%
\regfield{LEDC\_PARA\_UP\_LSCH\textcolor{red}{\it{n}}}{1}{4}{0}%
\regfield{LEDC\_IDLE\_LV\_LSCH\textcolor{red}{\it{n}}}{1}{3}{0}%
\regfield{LEDC\_SIG\_OUT\_EN\_LSCH\textcolor{red}{\it{n}}}{1}{2}{0}%
\regfield{LEDC\_TIMER\_SEL\_LSCH\textcolor{red}{\it{n}}}{2}{0}{{0}}%
\reglabel{Reset}\regnewline%
\begin{regdesc}\begin{reglist}
\item [LEDC\_PARA\_UP\_LSCH\textcolor{red}{\it{n}}]This bit is used to update register LEDC\_LSCH\textcolor{red}{\it{n}}\_HPOINT and LEDC\_LSCH\textcolor{red}{\it{n}}\_DUTY for low-speed channel \textcolor{red}{\it{n}}. (R/W)
\item [LEDC\_IDLE\_LV\_LSCH\textcolor{red}{\it{n}}]This bit is used to control the output value, when low-speed channel \textcolor{red}{\it{n}} is inactive. (R/W)
\item [LEDC\_SIG\_OUT\_EN\_LSCH\textcolor{red}{\it{n}}]This is the output enable control bit for low-speed channel \textcolor{red}{\it{n}}. (R/W)
\item [LEDC\_TIMER\_SEL\_LSCH\textcolor{red}{\it{n}}]There are four low-speed timers, the two bits are used to select one of them for low-speed channel \textcolor{red}{\it{n}}.  (R/W) \\0: select lstimer0; \\1: select lstimer1;  \\2: select lstimer2;  \\3: select lstimer3.
\end{reglist}\end{regdesc}
\end{register}


\begin{register}{H}{LEDC\_LSCH\textcolor{red}{\it{n}}\_HPOINT\_REG (\textcolor{red}{\it{n}}: 0-7)}{0xA4+0x14*\textcolor{red}{\it{n}}}\label{regdesc:LEDCLSCHnHPOINTREG}
\regfield{(reserved)}{12}{20}{{0x0000}}%
\regfield{LEDC\_HPOINT\_LSCH\textcolor{red}{\it{n}}}{20}{0}{{0x000000}}%
\reglabel{Reset}\regnewline%
\begin{regdesc}\begin{reglist}
\item [LEDC\_HPOINT\_LSCH\textcolor{red}{\it{n}}]The output value changes to high when lstimerx(x=[0,3]), selected by low-speed channel \textcolor{red}{\it{n}}, has reached LEDC\_HPOINT\_LSCH\textcolor{red}{\it{n}}[19:0]. (R/W)
\end{reglist}\end{regdesc}
\end{register}


\begin{register}{H}{LEDC\_LSCH\textcolor{red}{\it{n}}\_DUTY\_REG (\textcolor{red}{\it{n}}: 0-7)}{0xA8+0x14*\textcolor{red}{\it{n}}}\label{regdesc:LEDCLSCHnDUTYREG}
\regfield{(reserved)}{7}{25}{{0x00}}%
\regfield{LEDC\_DUTY\_LSCH\textcolor{red}{\it{n}}}{25}{0}{{0x0000000}}%
\reglabel{Reset}\regnewline%
\begin{regdesc}\begin{reglist}
\item [LEDC\_DUTY\_LSCH\textcolor{red}{\it{n}}]The register is used to control output duty. When lstimerx(x=[0,3]), chosen by low-speed channel \textcolor{red}{\it{n}}, has reached LEDC\_LPOINT\_LSCH\textcolor{red}{\it{n}},the output signal changes to low. (R/W) \\ LEDC\_LPOINT\_LSCH\textcolor{red}{\it{n}}=(LEDC\_HPOINT\_LSCH\textcolor{red}{\it{n}}[19:0]+LEDC\_DUTY\_LSCH\textcolor{red}{\it{n}}[24:4]) (1) \\LEDC\_LPOINT\_LSCH\textcolor{red}{\it{n}}=(LEDC\_HPOINT\_LSCH\textcolor{red}{\it{n}}[19:0]+LEDC\_DUTY\_LSCH\textcolor{red}{\it{n}}[24:4] +1) (2) \\    See the \hyperref[ledcfd:dithering]{Functional Description} for more information on when (1) or (2) is chosen.
\end{reglist}\end{regdesc}
\end{register}


\begin{register}{H}{LEDC\_LSCH\textcolor{red}{\it{n}}\_CONF1\_REG (\textcolor{red}{\it{n}}: 0-7)}{0xAC+0x14*\textcolor{red}{\it{n}}}\label{regdesc:LEDCLSCHnCONF1REG}
\regfield{LEDC\_DUTY\_START\_LSCH\textcolor{red}{\it{n}}}{1}{31}{0}%
\regfield{LEDC\_DUTY\_INC\_LSCH\textcolor{red}{\it{n}}}{1}{30}{1}%
\regfield{LEDC\_DUTY\_NUM\_LSCH\textcolor{red}{\it{n}}}{10}{20}{{0x000}}%
\regfield{LEDC\_DUTY\_CYCLE\_LSCH\textcolor{red}{\it{n}}}{10}{10}{{0x000}}%
\regfield{LEDC\_DUTY\_SCALE\_LSCH\textcolor{red}{\it{n}}}{10}{0}{{0x000}}%
\reglabel{Reset}\regnewline%
\begin{regdesc}\begin{reglist}
\item [LEDC\_DUTY\_START\_LSCH\textcolor{red}{\it{n}}]When LEDC\_DUTY\_NUM\_HSCH\textcolor{red}{\it{n}}, LEDC\_DUTY\_CYCLE\_HSCH\textcolor{red}{\it{n}} and LEDC\_DUTY\_SCALE\_HSCH\textcolor{red}{\it{n}} have been configured, these settings will not take effect until set LEDC\_DUTY\_START\_HSCH\textcolor{red}{\it{n}}. This bit is automatically cleared by hardware. (R/W)
\item [LEDC\_DUTY\_INC\_LSCH\textcolor{red}{\it{n}}]This register is used to increase or decrease the duty of output signal for low-speed channel \textcolor{red}{\it{n}}. (R/W)
\item [LEDC\_DUTY\_NUM\_LSCH\textcolor{red}{\it{n}}]This register is used to control the number of times the duty cycle is increased or decreased for low-speed channel \textcolor{red}{\it{n}}. (R/W)
\item [LEDC\_DUTY\_CYCLE\_LSCH\textcolor{red}{\it{n}}]This register is used to increase or decrease the duty every LEDC\_DUTY\_CYCLE\_LSCH\textcolor{red}{\it{n}} cycles for low-speed channel \textcolor{red}{\it{n}}. (R/W)
\item [LEDC\_DUTY\_SCALE\_LSCH\textcolor{red}{\it{n}}]This register is used to increase or decrease the step scale for low-speed channel \textcolor{red}{\it{n}}. (R/W)
\end{reglist}\end{regdesc}
\end{register}


\begin{register}{H}{LEDC\_LSCH\textcolor{red}{\it{n}}\_DUTY\_R\_REG (\textcolor{red}{\it{n}}: 0-7)}{0xB0+0x14*\textcolor{red}{\it{n}}}\label{regdesc:LEDCLSCHnDUTYRREG}
\regfield{(reserved)}{7}{25}{{0x00}}%
\regfield{LEDC\_DUTY\_LSCH\textcolor{red}{\it{n}}\_R}{25}{0}{{0x0000000}}%
\reglabel{Reset}\regnewline%
\begin{regdesc}\begin{reglist}
\item [LEDC\_DUTY\_LSCH\textcolor{red}{\it{n}}\_R]This register represents the current duty of the output signal for low-speed channel \textcolor{red}{\it{n}}. (RO)
\end{reglist}\end{regdesc}
\end{register}


\begin{register}{H}{LEDC\_HSTIMER\textcolor{red}{\it{x}}\_CONF\_REG (\textcolor{red}{\it{x}}: 0-3)}{0x140+8*\textcolor{red}{\it{x}}}\label{regdesc:LEDCHSTIMERnCONFREG}
\regfield{(reserved)}{6}{26}{{0x00}}%
\regfield{LEDC\_TICK\_SEL\_HSTIMER\textcolor{red}{\it{x}}}{1}{25}{0}%
\regfield{LEDC\_HSTIMER\textcolor{red}{\it{x}}\_RST}{1}{24}{1}%
\regfield{LEDC\_HSTIMER\textcolor{red}{\it{x}}\_PAUSE}{1}{23}{0}%
\regfield{LEDC\_CLK\_DIV\_NUM\_HSTIMER\textcolor{red}{\it{x}}}{18}{5}{{0x00000}}%
\regfield{LEDC\_HSTIMER\textcolor{red}{\it{x}}\_DUTY\_RES}{5}{0}{{0x00}}%
\reglabel{Reset}\regnewline%
\begin{regdesc}\begin{reglist}
\item [LEDC\_TICK\_SEL\_HSTIMER\textcolor{red}{\it{x}}]This bit is used to select APB\_CLK or REF\_TICK for high-speed timer \textcolor{red}{\it{x}}. (R/W) \\1: APB\_CLK;  \\0: REF\_TICK.
\item [LEDC\_HSTIMER\textcolor{red}{\it{x}}\_RST]This bit is used to reset high-speed timer \textcolor{red}{\it{x}}. The counter value will be 'zero' after reset. (R/W)
\item [LEDC\_HSTIMER\textcolor{red}{\it{x}}\_PAUSE]This bit is used to suspend the counter in high-speed timer \textcolor{red}{\it{x}}. (R/W)
\item [LEDC\_CLK\_DIV\_NUM\_HSTIMER\textcolor{red}{\it{x}}]This register is used to configure the division factor for the divider in high-speed timer \textcolor{red}{\it{x}}. The least significant eight bits represent the fractional part. (R/W)
\item [LEDC\_HSTIMER\textcolor{red}{\it{x}}\_DUTY\_RES]This register is used to control the range of the counter in high-speed timer \textcolor{red}{\it{x}}. The counter range is [0, 2$ ^{LEDC\_HSTIMER\textcolor{red}{\it{x}}\_DUTY\_RES}$], the maximum bit width for counter is 20. (R/W)
\end{reglist}\end{regdesc}
\end{register}


\begin{register}{H}{LEDC\_HSTIMER\textcolor{red}{\it{x}}\_VALUE\_REG (\textcolor{red}{\it{x}}: 0-3)}{0x144+8*\textcolor{red}{\it{x}}}\label{regdesc:LEDCHSTIMERnVALUEREG}
\regfield{(reserved)}{12}{20}{{0x0000}}%
\regfield{LEDC\_HSTIMER\textcolor{red}{\it{x}}\_CNT}{20}{0}{00000000000000000000}%
\reglabel{Reset}\regnewline%
\begin{regdesc}\begin{reglist}
\item [LEDC\_HSTIMER\textcolor{red}{\it{x}}\_CNT]Software can read this register to get the current counter value of high-speed timer \textcolor{red}{\it{x}}. (RO)
\end{reglist}\end{regdesc}
\end{register}


\begin{register}{H}{LEDC\_LSTIMER\textcolor{red}{\it{x}}\_CONF\_REG (\textcolor{red}{\it{x}}: 0-3)}{0x160+8*\textcolor{red}{\it{x}}}\label{regdesc:LEDCLSTIMERnCONFREG}
\regfield{(reserved)}{5}{27}{{0x00}}%
\regfield{LEDC\_LSTIMER\textcolor{red}{\it{x}}\_PARA\_UP}{1}{26}{0}%
\regfield{LEDC\_TICK\_SEL\_LSTIMER\textcolor{red}{\it{x}}}{1}{25}{0}%
\regfield{LEDC\_LSTIMER\textcolor{red}{\it{x}}\_RST}{1}{24}{1}%
\regfield{LEDC\_LSTIMER\textcolor{red}{\it{x}}\_PAUSE}{1}{23}{0}%
\regfield{LEDC\_CLK\_DIV\_NUM\_LSTIMER\textcolor{red}{\it{x}}}{18}{5}{{0x00000}}%
\regfield{LEDC\_LSTIMER\textcolor{red}{\it{x}}\_DUTY\_RES}{5}{0}{{0x00}}%
\reglabel{Reset}\regnewline%
\begin{regdesc}\begin{reglist}
\item [LEDC\_LSTIMER\textcolor{red}{\it{x}}\_PARA\_UP]Set this bit  to update  LEDC\_CLK\_DIV\_NUM\_LSTIME\textcolor{red}{\it{x}} and  LEDC\_LSTIMER\textcolor{red}{\it{x}}\_DUTY\_RES. (R/W)
\item [LEDC\_TICK\_SEL\_LSTIMER\textcolor{red}{\it{x}}]This bit is used to select RTC\_SLOW\_CLK or REF\_TICK for low-speed timer \textcolor{red}{\it{x}}. (R/W) \\1: RTC\_SLOW\_CLK; \\0: REF\_TICK.
\item [LEDC\_LSTIMER\textcolor{red}{\it{x}}\_RST]This bit is used to reset low-speed timer \textcolor{red}{\it{x}}. The counter will show 0 after reset. (R/W)
\item [LEDC\_LSTIMER\textcolor{red}{\it{x}}\_PAUSE]This bit is used to suspend the counter in low-speed timer \textcolor{red}{\it{x}}. (R/W)
\item [LEDC\_CLK\_DIV\_NUM\_LSTIMER\textcolor{red}{\it{x}}]This register is used to configure the division factor for the divider in low-speed timer \textcolor{red}{\it{x}}. The least significant eight bits represent the fractional part. (R/W)
\item [LEDC\_LSTIMER\textcolor{red}{\it{x}}\_DUTY\_RES]This register is used to control the range of the counter in low-speed timer \textcolor{red}{\it{x}}. The counter range is [0, 2$ ^{LEDC\_LSTIMER\textcolor{red}{\it{x}}\_DUTY\_RES}$], the max bit width for counter is 20. (R/W)
\end{reglist}\end{regdesc}
\end{register}


\begin{register}{H}{LEDC\_LSTIMER\textcolor{red}{\it{x}}\_VALUE\_REG (\textcolor{red}{\it{x}}: 0-3)}{0x164+8*\textcolor{red}{\it{x}}}\label{regdesc:LEDCLSTIMERnVALUEREG}
\regfield{(reserved)}{12}{20}{{0x0000}}%
\regfield{LEDC\_LSTIMER\textcolor{red}{\it{x}}\_CNT}{20}{0}{00000000000000000000}%
\reglabel{Reset}\regnewline%
\begin{regdesc}\begin{reglist}
\item [LEDC\_LSTIMER\textcolor{red}{\it{x}}\_CNT]Software can read this register to get the current counter value of low-speed timer \textcolor{red}{\it{x}}. (RO)
\end{reglist}\end{regdesc}
\end{register}


\begin{register}{H}{LEDC\_INT\_RAW\_REG}{0x0180}\label{regdesc:LEDCINTRAWREG}
\regfield{(reserved)}{8}{24}{00000000}%
\regfield{LEDC\_DUTY\_CHNG\_END\_LSCH7\_INT\_RAW}{1}{23}{0}%
\regfield{LEDC\_DUTY\_CHNG\_END\_LSCH6\_INT\_RAW}{1}{22}{0}%
\regfield{LEDC\_DUTY\_CHNG\_END\_LSCH5\_INT\_RAW}{1}{21}{0}%
\regfield{LEDC\_DUTY\_CHNG\_END\_LSCH4\_INT\_RAW}{1}{20}{0}%
\regfield{LEDC\_DUTY\_CHNG\_END\_LSCH3\_INT\_RAW}{1}{19}{0}%
\regfield{LEDC\_DUTY\_CHNG\_END\_LSCH2\_INT\_RAW}{1}{18}{0}%
\regfield{LEDC\_DUTY\_CHNG\_END\_LSCH1\_INT\_RAW}{1}{17}{0}%
\regfield{LEDC\_DUTY\_CHNG\_END\_LSCH0\_INT\_RAW}{1}{16}{0}%
\regfield{LEDC\_DUTY\_CHNG\_END\_HSCH7\_INT\_RAW}{1}{15}{0}%
\regfield{LEDC\_DUTY\_CHNG\_END\_HSCH6\_INT\_RAW}{1}{14}{0}%
\regfield{LEDC\_DUTY\_CHNG\_END\_HSCH5\_INT\_RAW}{1}{13}{0}%
\regfield{LEDC\_DUTY\_CHNG\_END\_HSCH4\_INT\_RAW}{1}{12}{0}%
\regfield{LEDC\_DUTY\_CHNG\_END\_HSCH3\_INT\_RAW}{1}{11}{0}%
\regfield{LEDC\_DUTY\_CHNG\_END\_HSCH2\_INT\_RAW}{1}{10}{0}%
\regfield{LEDC\_DUTY\_CHNG\_END\_HSCH1\_INT\_RAW}{1}{9}{0}%
\regfield{LEDC\_DUTY\_CHNG\_END\_HSCH0\_INT\_RAW}{1}{8}{0}%
\regfield{LEDC\_LSTIMER3\_OVF\_INT\_RAW}{1}{7}{0}%
\regfield{LEDC\_LSTIMER2\_OVF\_INT\_RAW}{1}{6}{0}%
\regfield{LEDC\_LSTIMER1\_OVF\_INT\_RAW}{1}{5}{0}%
\regfield{LEDC\_LSTIMER0\_OVF\_INT\_RAW}{1}{4}{0}%
\regfield{LEDC\_HSTIMER3\_OVF\_INT\_RAW}{1}{3}{0}%
\regfield{LEDC\_HSTIMER2\_OVF\_INT\_RAW}{1}{2}{0}%
\regfield{LEDC\_HSTIMER1\_OVF\_INT\_RAW}{1}{1}{0}%
\regfield{LEDC\_HSTIMER0\_OVF\_INT\_RAW}{1}{0}{0}%
\reglabel{Reset}\regnewline%
\begin{regdesc}\begin{reglist}
\item [LEDC\_DUTY\_CHNG\_END\_LSCH\textcolor{red}{\it{n}}\_INT\_RAW]The raw interrupt status bit for the \hyperref[int:LEDCDUTYCHNGENDLSCHnINT]{LEDC\_DUTY\_CHNG\_END\_LSCH\textcolor{red}{\it{n}}\_INT} interrupt. (RO)
\item [LEDC\_DUTY\_CHNG\_END\_HSCH\textcolor{red}{\it{n}}\_INT\_RAW]The raw interrupt status bit for the \hyperref[int:LEDCDUTYCHNGENDHSCHnINT]{LEDC\_DUTY\_CHNG\_END\_HSCH\textcolor{red}{\it{n}}\_INT} interrupt. (RO)
\item [LEDC\_LSTIMER\textcolor{red}{\it{x}}\_OVF\_INT\_RAW]The raw interrupt status bit for the \hyperref[int:LEDCLSTIMERxOVFINT]{LEDC\_LSTIMER\textcolor{red}{\it{x}}\_OVF\_INT} interrupt. (RO)
\item [LEDC\_HSTIMER\textcolor{red}{\it{x}}\_OVF\_INT\_RAW]The raw interrupt status bit for the \hyperref[int:LEDCHSTIMERxOVFINT]{LEDC\_HSTIMER\textcolor{red}{\it{x}}\_OVF\_INT} interrupt. (RO)
\end{reglist}\end{regdesc}
\end{register}


\begin{register}{H}{LEDC\_INT\_ST\_REG}{0x0184}\label{regdesc:LEDCINTSTREG}
\regfield{(reserved)}{8}{24}{00000000}%
\regfield{LEDC\_DUTY\_CHNG\_END\_LSCH7\_INT\_ST}{1}{23}{0}%
\regfield{LEDC\_DUTY\_CHNG\_END\_LSCH6\_INT\_ST}{1}{22}{0}%
\regfield{LEDC\_DUTY\_CHNG\_END\_LSCH5\_INT\_ST}{1}{21}{0}%
\regfield{LEDC\_DUTY\_CHNG\_END\_LSCH4\_INT\_ST}{1}{20}{0}%
\regfield{LEDC\_DUTY\_CHNG\_END\_LSCH3\_INT\_ST}{1}{19}{0}%
\regfield{LEDC\_DUTY\_CHNG\_END\_LSCH2\_INT\_ST}{1}{18}{0}%
\regfield{LEDC\_DUTY\_CHNG\_END\_LSCH1\_INT\_ST}{1}{17}{0}%
\regfield{LEDC\_DUTY\_CHNG\_END\_LSCH0\_INT\_ST}{1}{16}{0}%
\regfield{LEDC\_DUTY\_CHNG\_END\_HSCH7\_INT\_ST}{1}{15}{0}%
\regfield{LEDC\_DUTY\_CHNG\_END\_HSCH6\_INT\_ST}{1}{14}{0}%
\regfield{LEDC\_DUTY\_CHNG\_END\_HSCH5\_INT\_ST}{1}{13}{0}%
\regfield{LEDC\_DUTY\_CHNG\_END\_HSCH4\_INT\_ST}{1}{12}{0}%
\regfield{LEDC\_DUTY\_CHNG\_END\_HSCH3\_INT\_ST}{1}{11}{0}%
\regfield{LEDC\_DUTY\_CHNG\_END\_HSCH2\_INT\_ST}{1}{10}{0}%
\regfield{LEDC\_DUTY\_CHNG\_END\_HSCH1\_INT\_ST}{1}{9}{0}%
\regfield{LEDC\_DUTY\_CHNG\_END\_HSCH0\_INT\_ST}{1}{8}{0}%
\regfield{LEDC\_LSTIMER3\_OVF\_INT\_ST}{1}{7}{0}%
\regfield{LEDC\_LSTIMER2\_OVF\_INT\_ST}{1}{6}{0}%
\regfield{LEDC\_LSTIMER1\_OVF\_INT\_ST}{1}{5}{0}%
\regfield{LEDC\_LSTIMER0\_OVF\_INT\_ST}{1}{4}{0}%
\regfield{LEDC\_HSTIMER3\_OVF\_INT\_ST}{1}{3}{0}%
\regfield{LEDC\_HSTIMER2\_OVF\_INT\_ST}{1}{2}{0}%
\regfield{LEDC\_HSTIMER1\_OVF\_INT\_ST}{1}{1}{0}%
\regfield{LEDC\_HSTIMER0\_OVF\_INT\_ST}{1}{0}{0}%
\reglabel{Reset}\regnewline%
\begin{regdesc}\begin{reglist}
\item [LEDC\_DUTY\_CHNG\_END\_LSCH\textcolor{red}{\it{n}}\_INT\_ST]The masked interrupt status bit for the \hyperref[int:LEDCDUTYCHNGENDLSCHnINT]{LEDC\_DUTY\_CHNG\_END\_LSCH\textcolor{red}{\it{n}}\_INT} interrupt. (RO)
\item [LEDC\_DUTY\_CHNG\_END\_HSCH\textcolor{red}{\it{n}}\_INT\_ST]The masked interrupt status bit for the \hyperref[int:LEDCDUTYCHNGENDHSCHnINT]{LEDC\_DUTY\_CHNG\_END\_HSCH\textcolor{red}{\it{n}}\_INT} interrupt. (RO)
\item [LEDC\_LSTIMER\textcolor{red}{\it{x}}\_OVF\_INT\_ST]The masked interrupt status bit for the \hyperref[int:LEDCLSTIMERxOVFINT]{LEDC\_LSTIMER\textcolor{red}{\it{x}}\_OVF\_INT} interrupt. (RO)
\item [LEDC\_HSTIMER\textcolor{red}{\it{x}}\_OVF\_INT\_ST]The masked interrupt status bit for the \hyperref[int:LEDCHSTIMERxOVFINT]{LEDC\_HSTIMER\textcolor{red}{\it{x}}\_OVF\_INT} interrupt. (RO)
\end{reglist}\end{regdesc}
\end{register}


\begin{register}{H}{LEDC\_INT\_ENA\_REG}{0x0188}\label{regdesc:LEDCINTENAREG}
\regfield{(reserved)}{8}{24}{00000000}%
\regfield{LEDC\_DUTY\_CHNG\_END\_LSCH7\_INT\_ENA}{1}{23}{0}%
\regfield{LEDC\_DUTY\_CHNG\_END\_LSCH6\_INT\_ENA}{1}{22}{0}%
\regfield{LEDC\_DUTY\_CHNG\_END\_LSCH5\_INT\_ENA}{1}{21}{0}%
\regfield{LEDC\_DUTY\_CHNG\_END\_LSCH4\_INT\_ENA}{1}{20}{0}%
\regfield{LEDC\_DUTY\_CHNG\_END\_LSCH3\_INT\_ENA}{1}{19}{0}%
\regfield{LEDC\_DUTY\_CHNG\_END\_LSCH2\_INT\_ENA}{1}{18}{0}%
\regfield{LEDC\_DUTY\_CHNG\_END\_LSCH1\_INT\_ENA}{1}{17}{0}%
\regfield{LEDC\_DUTY\_CHNG\_END\_LSCH0\_INT\_ENA}{1}{16}{0}%
\regfield{LEDC\_DUTY\_CHNG\_END\_HSCH7\_INT\_ENA}{1}{15}{0}%
\regfield{LEDC\_DUTY\_CHNG\_END\_HSCH6\_INT\_ENA}{1}{14}{0}%
\regfield{LEDC\_DUTY\_CHNG\_END\_HSCH5\_INT\_ENA}{1}{13}{0}%
\regfield{LEDC\_DUTY\_CHNG\_END\_HSCH4\_INT\_ENA}{1}{12}{0}%
\regfield{LEDC\_DUTY\_CHNG\_END\_HSCH3\_INT\_ENA}{1}{11}{0}%
\regfield{LEDC\_DUTY\_CHNG\_END\_HSCH2\_INT\_ENA}{1}{10}{0}%
\regfield{LEDC\_DUTY\_CHNG\_END\_HSCH1\_INT\_ENA}{1}{9}{0}%
\regfield{LEDC\_DUTY\_CHNG\_END\_HSCH0\_INT\_ENA}{1}{8}{0}%
\regfield{LEDC\_LSTIMER3\_OVF\_INT\_ENA}{1}{7}{0}%
\regfield{LEDC\_LSTIMER2\_OVF\_INT\_ENA}{1}{6}{0}%
\regfield{LEDC\_LSTIMER1\_OVF\_INT\_ENA}{1}{5}{0}%
\regfield{LEDC\_LSTIMER0\_OVF\_INT\_ENA}{1}{4}{0}%
\regfield{LEDC\_HSTIMER3\_OVF\_INT\_ENA}{1}{3}{0}%
\regfield{LEDC\_HSTIMER2\_OVF\_INT\_ENA}{1}{2}{0}%
\regfield{LEDC\_HSTIMER1\_OVF\_INT\_ENA}{1}{1}{0}%
\regfield{LEDC\_HSTIMER0\_OVF\_INT\_ENA}{1}{0}{0}%
\reglabel{Reset}\regnewline%
\begin{regdesc}\begin{reglist}
\item [LEDC\_DUTY\_CHNG\_END\_LSCH\textcolor{red}{\it{n}}\_INT\_ENA]The interrupt enable bit for the \hyperref[int:LEDCDUTYCHNGENDLSCHnINT]{LEDC\_DUTY\_CHNG\_END\_LSCH\textcolor{red}{\it{n}}\_INT} interrupt. (R/W)
\item [LEDC\_DUTY\_CHNG\_END\_HSCH\textcolor{red}{\it{n}}\_INT\_ENA]The interrupt enable bit for the \hyperref[int:LEDCDUTYCHNGENDHSCHnINT]{LEDC\_DUTY\_CHNG\_END\_HSCH\textcolor{red}{\it{n}}\_INT} interrupt. (R/W)
\item [LEDC\_LSTIMER\textcolor{red}{\it{x}}\_OVF\_INT\_ENA]The interrupt enable bit for the \hyperref[int:LEDCLSTIMERxOVFINT]{LEDC\_LSTIMER\textcolor{red}{\it{x}}\_OVF\_INT} interrupt. (R/W)
\item [LEDC\_HSTIMER\textit{\textcolor{red}{\it{x}}}\_OVF\_INT\_ENA]The interrupt enable bit for the \hyperref[int:LEDCHSTIMERxOVFINT]{LEDC\_HSTIMER\textcolor{red}{\it{x}}\_OVF\_INT} interrupt. (R/W)
\end{reglist}\end{regdesc}
\end{register}


\begin{register}{H}{LEDC\_INT\_CLR\_REG}{0x018C}\label{regdesc:LEDCINTCLRREG}
\regfield{(reserved)}{8}{24}{00000000}%
\regfield{LEDC\_DUTY\_CHNG\_END\_LSCH7\_INT\_CLR}{1}{23}{0}%
\regfield{LEDC\_DUTY\_CHNG\_END\_LSCH6\_INT\_CLR}{1}{22}{0}%
\regfield{LEDC\_DUTY\_CHNG\_END\_LSCH5\_INT\_CLR}{1}{21}{0}%
\regfield{LEDC\_DUTY\_CHNG\_END\_LSCH4\_INT\_CLR}{1}{20}{0}%
\regfield{LEDC\_DUTY\_CHNG\_END\_LSCH3\_INT\_CLR}{1}{19}{0}%
\regfield{LEDC\_DUTY\_CHNG\_END\_LSCH2\_INT\_CLR}{1}{18}{0}%
\regfield{LEDC\_DUTY\_CHNG\_END\_LSCH1\_INT\_CLR}{1}{17}{0}%
\regfield{LEDC\_DUTY\_CHNG\_END\_LSCH0\_INT\_CLR}{1}{16}{0}%
\regfield{LEDC\_DUTY\_CHNG\_END\_HSCH7\_INT\_CLR}{1}{15}{0}%
\regfield{LEDC\_DUTY\_CHNG\_END\_HSCH6\_INT\_CLR}{1}{14}{0}%
\regfield{LEDC\_DUTY\_CHNG\_END\_HSCH5\_INT\_CLR}{1}{13}{0}%
\regfield{LEDC\_DUTY\_CHNG\_END\_HSCH4\_INT\_CLR}{1}{12}{0}%
\regfield{LEDC\_DUTY\_CHNG\_END\_HSCH3\_INT\_CLR}{1}{11}{0}%
\regfield{LEDC\_DUTY\_CHNG\_END\_HSCH2\_INT\_CLR}{1}{10}{0}%
\regfield{LEDC\_DUTY\_CHNG\_END\_HSCH1\_INT\_CLR}{1}{9}{0}%
\regfield{LEDC\_DUTY\_CHNG\_END\_HSCH0\_INT\_CLR}{1}{8}{0}%
\regfield{LEDC\_LSTIMER3\_OVF\_INT\_CLR}{1}{7}{0}%
\regfield{LEDC\_LSTIMER2\_OVF\_INT\_CLR}{1}{6}{0}%
\regfield{LEDC\_LSTIMER1\_OVF\_INT\_CLR}{1}{5}{0}%
\regfield{LEDC\_LSTIMER0\_OVF\_INT\_CLR}{1}{4}{0}%
\regfield{LEDC\_HSTIMER3\_OVF\_INT\_CLR}{1}{3}{0}%
\regfield{LEDC\_HSTIMER2\_OVF\_INT\_CLR}{1}{2}{0}%
\regfield{LEDC\_HSTIMER1\_OVF\_INT\_CLR}{1}{1}{0}%
\regfield{LEDC\_HSTIMER0\_OVF\_INT\_CLR}{1}{0}{0}%
\reglabel{Reset}\regnewline%
\begin{regdesc}\begin{reglist}
\item [LEDC\_DUTY\_CHNG\_END\_LSCH\textcolor{red}{\it{n}}\_INT\_CLR]Set this bit to clear the \hyperref[int:LEDCDUTYCHNGENDLSCHnINT]{LEDC\_DUTY\_CHNG\_END\_LSCH\textcolor{red}{\it{n}}\_INT} interrupt. (WO)
\item [LEDC\_DUTY\_CHNG\_END\_HSCH\textcolor{red}{\it{n}}\_INT\_CLR]Set this bit to clear the \hyperref[int:LEDCDUTYCHNGENDHSCHnINT]{LEDC\_DUTY\_CHNG\_END\_HSCH\textcolor{red}{\it{n}}\_INT} interrupt. (WO)
\item [LEDC\_LSTIMER\textcolor{red}{\it{x}}\_OVF\_INT\_CLR]Set this bit to clear the \hyperref[int:LEDCLSTIMERxOVFINT]{LEDC\_LSTIMER\textcolor{red}{\it{x}}\_OVF\_INT} interrupt. (WO)
\item [LEDC\_HSTIMER\textcolor{red}{\it{x}}\_OVF\_INT\_CLR]Set this bit to clear the \hyperref[int:LEDCHSTIMERxOVFINT]{LEDC\_HSTIMER\textcolor{red}{\it{x}}\_OVF\_INT} interrupt. (WO)
\end{reglist}\end{regdesc}
\end{register}


\begin{register}{H}{LEDC\_CONF\_REG}{0x0190}\label{regdesc:LEDCCONFREG}
\regfield{(reserved)}{31}{1}{0000000000000000000000000000000}%
\regfield{LEDC\_APB\_CLK\_SEL}{1}{0}{0}%
\reglabel{Reset}\regnewline%
\begin{regdesc}\begin{reglist}
\item [LEDC\_APB\_CLK\_SEL]This bit is used to set the frequency of RTC\_SLOW\_CLK. (R/W) \\
0: 8 MHz; \\
1: 80 MHz.
\end{reglist}\end{regdesc}
\end{register}
